% GENERATED FILE. Edit canonical lessons, then run npm run generate:books.
% canonical-source-hash: fnv1a64:9fe6cec542f15964
% canonical-lessons: LA-C381-causidicus, LA-C381-cliens, LA-C381-navicularius, LA-C381-librarius, LA-C381-agrimensor

\chapter{Pleaders, Clients and Clerks}
\label{ch:la-pleaders-clients-and-clerks}
\hlchaptermodality{381}

\begin{chapteropening}
\textbf{By the end of this chapter:} \emph{I can name an advocate and his client, a shipowner, a copyist and a land surveyor.}

Complete the last lesson of chapter 381: I can name an advocate and his client, a shipowner, a copyist and a land surveyor.
\end{chapteropening}

\section[causidicus, causidicī]{causidicus, causidicī — an advocate, a pleader for hire}

\label{lesson:LA-C381-causidicus}



\begin{quote}
Before the new one: say the Latin for to do with energy, then the Latin for to tire out.
\end{quote}

\subsection*{You'll want to know: causidicus, causidicī}
\textbf{causidicus, causidicī} — \textquotedblleft{}an advocate, a pleader for hire\textquotedblright{}.

\begin{quote}
Causidicus prō reō dīxit.
\end{quote}

\begin{quote}
\emph{The advocate spoke for the accused.}
\end{quote}

\emph{causidicus, causidicī} m., from \emph{causa}, a case, and \emph{dīcō}: a case-speaker, often said with a little scorn.

\begin{cousinweb}
\emph{causidicus} is built from \emph{causa} \textquotedblleft{}a case at law\textquotedblright{} and \emph{dīcō} \textquotedblleft{}to say\textquotedblright{} (Lewis \& Short): one who speaks cases.
\end{cousinweb}

\begin{grammarlens}[title={What you've built}]
One more word for telling, explaining and describing at B1.
\end{grammarlens}

\subsection*{Your turn}
\begin{itemize}
\raggedright
  \item \emph{Say it:} \emph{causidicus, causidicī}
  \item \emph{Say it:} \emph{causidicus, causidicī}, once more
  \item \emph{Say it:} a sentence of your own with \emph{causidicus, causidicī}
  \item \emph{Recall:} say the Latin for to do with energy, then the Latin for to tire out, then say \emph{causidicus, causidicī} again
\end{itemize}

\begin{tcolorbox}[breakable,colback=teal!4,colframe=teal!35!black,title={Before you move on}]
What does causidicus, causidicī mean? (\textquotedblleft{}An advocate, a pleader for hire\textquotedblright{}.) Say it once more.
\end{tcolorbox}
\section[cliēns, clientis]{cliēns, clientis — a client, a dependant}

\label{lesson:LA-C381-cliens}



\begin{quote}
Before the new one: say the Latin for to tire out, then the Latin for an advocate.
\end{quote}

\subsection*{You'll want to know: cliēns, clientis}
\textbf{cliēns, clientis} — \textquotedblleft{}a client, a dependant\textquotedblright{}.

\begin{quote}
Clientēs māne ad domum patrōnī veniunt.
\end{quote}

\begin{quote}
\emph{The clients come to their patron's house in the morning.}
\end{quote}

\emph{cliēns, clientis} m. A client owed his patron greeting and support; the patron owed him protection and help at law. English \emph{client}.

\begin{cousinweb}
\emph{cliēns} is usually traced to an old verb \emph{clueō} \textquotedblleft{}to hear, to be called\textquotedblright{} (Lewis \& Short): one who listens to his patron. English \emph{client} comes from it.
\end{cousinweb}

\begin{grammarlens}[title={What you've built}]
One more word for telling, explaining and describing at B1.
\end{grammarlens}

\subsection*{Your turn}
\begin{itemize}
\raggedright
  \item \emph{Say it:} \emph{cliēns, clientis}
  \item \emph{Say it:} \emph{cliēns, clientis}, once more
  \item \emph{Say it:} a sentence of your own with \emph{cliēns, clientis}
  \item \emph{Recall:} say the Latin for to tire out, then the Latin for an advocate, then say \emph{cliēns, clientis} again
\end{itemize}

\begin{tcolorbox}[breakable,colback=teal!4,colframe=teal!35!black,title={Before you move on}]
What does cliēns, clientis mean? (\textquotedblleft{}A client, a dependant\textquotedblright{}.) Say it once more.
\end{tcolorbox}
\section[nāviculārius, nāviculāriī]{nāviculārius, nāviculāriī — a shipowner, a shipper}

\label{lesson:LA-C381-navicularius}



\begin{quote}
Before the new one: say the Latin for an advocate, then the Latin for a client.
\end{quote}

\subsection*{You'll want to know: nāviculārius, nāviculāriī}
\textbf{nāviculārius, nāviculāriī} — \textquotedblleft{}a shipowner, a shipper\textquotedblright{}.

\begin{quote}
Nāviculārius frūmentum trāns mare vehit.
\end{quote}

\begin{quote}
\emph{The shipowner carries grain across the sea.}
\end{quote}

\emph{nāviculārius, nāviculāriī} m., built on a little word for a ship, from \emph{nāvis}. The great grain ships of Rome were run by \emph{nāviculāriī}.

\begin{cousinweb}
\emph{nāviculārius} is built on \emph{nāvicula} \textquotedblleft{}a little ship\textquotedblright{} (Lewis \& Short), a small form of \emph{nāvis} \textquotedblleft{}a ship\textquotedblright{}.
\end{cousinweb}

\begin{grammarlens}[title={What you've built}]
One more word for telling, explaining and describing at B1.
\end{grammarlens}

\subsection*{Your turn}
\begin{itemize}
\raggedright
  \item \emph{Say it:} \emph{nāviculārius, nāviculāriī}
  \item \emph{Say it:} \emph{nāviculārius, nāviculāriī}, once more
  \item \emph{Say it:} a sentence of your own with \emph{nāviculārius, nāviculāriī}
  \item \emph{Recall:} say the Latin for an advocate, then the Latin for a client, then say \emph{nāviculārius, nāviculāriī} again
\end{itemize}

\begin{tcolorbox}[breakable,colback=teal!4,colframe=teal!35!black,title={Before you move on}]
What does nāviculārius, nāviculāriī mean? (\textquotedblleft{}A shipowner, a shipper\textquotedblright{}.) Say it once more.
\end{tcolorbox}
\section[librārius, librāriī]{librārius, librāriī — a copyist, a scribe; a bookseller}

\label{lesson:LA-C381-librarius}



\begin{quote}
Before the new one: say the Latin for a client, then the Latin for a shipowner.
\end{quote}

\subsection*{You'll want to know: librārius, librāriī}
\textbf{librārius, librāriī} — \textquotedblleft{}a copyist, a scribe; a bookseller\textquotedblright{}.

\begin{quote}
Librārius librōs dīligenter scrībit.
\end{quote}

\begin{quote}
\emph{The copyist writes out books carefully.}
\end{quote}

\emph{librārius, librāriī} m., built on \emph{liber}, a book, with a short \emph{i}. Copyists were often slaves trained to write a fast clear hand.

\begin{cousinweb}
\emph{librārius} is built on \emph{liber, librī} \textquotedblleft{}a book\textquotedblright{} (Lewis \& Short), with the ending \emph{-ārius} of a man busy with something. English \emph{library} comes from the same word through French.
\end{cousinweb}

\begin{grammarlens}[title={What you've built}]
One more word for telling, explaining and describing at B1.
\end{grammarlens}

\subsection*{Your turn}
\begin{itemize}
\raggedright
  \item \emph{Say it:} \emph{librārius, librāriī}
  \item \emph{Say it:} \emph{librārius, librāriī}, once more
  \item \emph{Say it:} a sentence of your own with \emph{librārius, librāriī}
  \item \emph{Recall:} say the Latin for a client, then the Latin for a shipowner, then say \emph{librārius, librāriī} again
\end{itemize}

\begin{tcolorbox}[breakable,colback=teal!4,colframe=teal!35!black,title={Before you move on}]
What does librārius, librāriī mean? (\textquotedblleft{}A copyist, a scribe; a bookseller\textquotedblright{}.) Say it once more.
\end{tcolorbox}
\section[agrīmēnsor, agrīmēnsōris]{agrīmēnsor, agrīmēnsōris — a land surveyor}

\label{lesson:LA-C381-agrimensor}



\begin{quote}
Before the new one: say the Latin for a shipowner, then the Latin for a copyist.
\end{quote}

\subsection*{You'll want to know: agrīmēnsor, agrīmēnsōris}
\textbf{agrīmēnsor, agrīmēnsōris} — \textquotedblleft{}a land surveyor\textquotedblright{}.

\begin{quote}
Agrīmēnsor fīnēs agrī mētītur.
\end{quote}

\begin{quote}
\emph{The surveyor measures the boundaries of the field.}
\end{quote}

\emph{agrīmēnsor, agrīmēnsōris} m., from \emph{agrī}, of a field, and \emph{mētior}, to measure (participle \emph{mēnsus}). A late word, in Ammianus and Cassiodorus.

\begin{cousinweb}
\emph{agrīmēnsor} is built from \emph{agrī}, the genitive of \emph{ager} \textquotedblleft{}a field\textquotedblright{} and \emph{mēnsor} \textquotedblleft{}a measurer\textquotedblright{} (Lewis \& Short), the agent noun of \emph{mētior} \textquotedblleft{}to measure\textquotedblright{}.
\end{cousinweb}

\begin{grammarlens}[title={What you've built}]
One more word for telling, explaining and describing at B1.
\end{grammarlens}

\subsection*{Your turn}
\begin{itemize}
\raggedright
  \item \emph{Say it:} \emph{agrīmēnsor, agrīmēnsōris}
  \item \emph{Say it:} \emph{agrīmēnsor, agrīmēnsōris}, once more
  \item \emph{Say it:} a sentence of your own with \emph{agrīmēnsor, agrīmēnsōris}
  \item \emph{Recall:} say the Latin for a shipowner, then the Latin for a copyist, then say \emph{agrīmēnsor, agrīmēnsōris} again
\end{itemize}

\begin{tcolorbox}[breakable,colback=teal!4,colframe=teal!35!black,title={Before you move on}]
What does agrīmēnsor, agrīmēnsōris mean? (\textquotedblleft{}A land surveyor\textquotedblright{}.) Say it once more.
\end{tcolorbox}
